%-------------------------------------------------------------------------------
% SECTION TITLE
%-------------------------------------------------------------------------------
\cvsection{项目经历}


%-------------------------------------------------------------------------------
% CONTENT
%-------------------------------------------------------------------------------
\begin{cventries}

%---------------------------------------------------------
  \cventry
    {个人项目}
    {A Simplified vLLM Core}
    {LLM 推理系统}
    {2025年12月 -- 2026年5月}
    {
      \begin{cvitems}
        \item {基于 PyTorch/CUDA 从零实现简化版 vLLM 推理核心，完成 Engine、Scheduler、BlockManager 等模块搭建，打通 LLM 请求调度与执行链路。}
        \item {实现 Paged KV Cache、Continuous Batching 与 prefill/decode 混合调度，支持多请求离线并发推理。}
        \item {实现 AWQ 量化，降低模型权重存储开销。}
      \end{cvitems}
    }

%---------------------------------------------------------
  \cventry
    {\textcolor{awesome-red}{毕业设计} \textbf{\textcolor{awesome-red}{（进行中）}}}
    {Possible Optimizations for MoE Inference in vLLM}
    {米兰理工大学}
    {2026年5月 -- 至今}
    {
      \begin{cvitems}
        \item {研究 Scheduling For Experts Movement：在 EPLB 中发生 Expert 迁移时，在不改变算法结果的前提下，对同一批 Batch 重新组合以缓解 Hotspot；目前已完成 ILP 建模并证明可得到更优组合，后续将在真实 GPU 集群上测试。}
        \item {研究在 EPLB 迁移阶段降低 Expert 迁移优先级，以提升推理吞吐、尾延迟稳定性与在线服务连续性；这一 Feature 的主要难点不在代码实现，而在实验设计与 A/B 对齐。}
      \end{cvitems}
    }

%---------------------------------------------------------
\end{cventries}
